\subsubsection{Unitaries from Hamiltonians}\label{sec:unitaries-from-hamiltonians}

Until now, we've talked about Hamiltonians: Cost Hamiltonian $H_C$ and Mixer Hamiltonian $H_M$. But a \textbf{quantum computer cannot execute Hamiltonians}. It executes \textbf{quantum gates} (unitary operators). So the obvious question is: \emph{\textbf{how do we transform a Hamiltonian into something that a quantum computer can execute?}} The answer is the bridge between \textbf{physics (Hamiltonians)} and \textbf{quantum circuits (gates)}.

\highspace
\textcolor{Red2}{\faIcon{exclamation-triangle} \textbf{The problem.}} Suppose we have built the Cost Hamiltonian $H_C$ for a given problem. Great! \emph{But now what?} Can we execute $H_C$ on IBM Quantum computers? No, we cannot. Quantum hardware understands only \textbf{unitary operators} (quantum gates). So we need a translation from Hamiltonians to unitary operators. The question is: \emph{how do we do that?}

\highspace
\textcolor{Green3}{\faIcon{check-circle} \textbf{The fundamental theorem}} \begin{theorem}
    A unitary operator $U$ can be represented as the exponential of a Hamiltonian $H$:
    \begin{equation}\label{eq:unitary-from-hamiltonian}
        U = e^{-i \alpha H} = \exp\left(-i \alpha H\right)
    \end{equation}
\end{theorem}
\noindent
This equation comes from \textbf{quantum mechanics}. Remember that the time evolution of a quantum system is given by the Schrödinger equation:
\begin{equation*}
    i \hbar \frac{\partial}{\partial t} \ket{\psi(t)} = H \ket{\psi(t)}
\end{equation*}
The solution of this equation is (assuming $H$ is time-independent):
\begin{equation*}
    \ket{\psi(t)} = e^{-i H t / \hbar} \ket{\psi(0)}
\end{equation*}
Where the operator $U(t) = e^{-i H t / \hbar}$ is exactly the \textbf{time-evolution operator}, which tells us how the quantum state evolves after a physical time $t$. 

\highspace
Unlike VQE, where physical time $t$ is considered, in QAOA, the Hamiltonian $H_C$ does not describe a physical system evolving naturally. Rather, it is an artificial Hamiltonian that encodes the objective function of an optimization problem. We still call it ``\emph{time-evolution}'' because, mathematically, it is the same operation. For this reason, the time-evolution operator can be generalized to the QAOA case, since we are \textbf{not concerned with physical time and only need a unitary generated by a Hamiltonian}. Thus, we can write:
\begin{equation*}
    U(t) = e^{-i H t / \hbar} \quad \Rightarrow \quad U = e^{-i \alpha H}, \quad \alpha = \dfrac{t}{\hbar}
\end{equation*}
\textcolor{Green3}{\faIcon{question-circle} \textbf{Why introduce $\alpha$?}} Because $\alpha$ becomes a \textbf{free parameter}. So instead of saying ``\emph{let the system evolve for $t = 0.73$ ns}'', we say ``\emph{apply the unitary with parameter $\alpha = 0.37$}'', since $U$ may represent a quantum gate, a variational layer, a QAOA cost operator or a simulation step. In other words, $\alpha$ is a \textbf{scalar variational parameter} that we can optimize to find the best solution to our problem.

\highspace
The $\alpha$ parameter controls the \textbf{strength of the unitary operator}. For example, if $\alpha = 0$, then $U = e^{-i 0 H} = I$, which means that the unitary does nothing. If $\alpha$ is very large, then $U$ will have a strong effect on the quantum state, causing a noticeable evolution.

\highspace
\begin{examplebox}[: Connection with rotations]
    The unitary operator $U = e^{-i \alpha H}$ can be interpreted as a rotation in the Hilbert space. The parameter $\alpha$ controls the angle of rotation, and the Hamiltonian $H$ determines the axis of rotation.

    \highspace
    For example, we can retrieve all the Pauli rotations (\autopageref{eq:rotation-operator-general}) from the general formula $U = e^{-i \alpha H}$ by choosing the appropriate Hamiltonian $H$:
    \begin{itemize}
        \item Pauli-X rotation:
        \begin{equation*}
            R_X(\theta) = e^{-i \frac{\theta}{2} X} = \exp\left(-i \frac{\theta}{2} \begin{pmatrix}
                0 & 1 \\
                1 & 0
            \end{pmatrix}\right)
        \end{equation*}
        \item Pauli-Y rotation:
        \begin{equation*}
            R_Y(\theta) = e^{-i \frac{\theta}{2} Y} = \exp\left(-i \frac{\theta}{2} \begin{pmatrix}
                0 & -i \\
                i & 0
            \end{pmatrix}\right)
        \end{equation*}
        \item Pauli-Z rotation:
        \begin{equation*}
            R_Z(\theta) = e^{-i \frac{\theta}{2} Z} = \exp\left(-i \frac{\theta}{2} \begin{pmatrix}
                1 & 0 \\
                0 & -1
            \end{pmatrix}\right)
        \end{equation*}
    \end{itemize}
\end{examplebox}